  \documentclass{article}
\usepackage{amsmath}
\usepackage{graphicx}
\graphicspath{ {images/} }
\usepackage[spanish]{babel}
\usepackage{bbm}
\usepackage{amsthm}
\usepackage{authblk}
\usepackage{hyperref}
\usepackage[utf8]{inputenc}
\usepackage{mathrsfs}
\usepackage{amsfonts}
\usepackage{amssymb}
\usepackage{enumerate}
\usepackage[left=3cm,right=2cm,top=2cm,bottom=2cm]{geometry}
\usepackage{epsfig}
 \renewcommand{\baselinestretch}{0.93}
 \thispagestyle{empty}
 \pagestyle{empty}
\setlength{\textheight}{53pc} \setlength{\textwidth} {36pc}
\setlength{\topmargin}{-4mm}
\usepackage{multicol}
\newcommand*{\email}[1]{%
    \normalsize\href{mailto:#1}{#1}\par
    }
\setlength{\columnsep}{1cm}
\newcommand{\N}{\mathbb{N}}
\newcommand{\calM}{\cal{M}}
\newcommand{\calP}{\cal{P}}
\newcommand{\F}{\mathbb{F}}
\newcommand{\R}{\mathbb{R}}
\newcommand{\Z}{\mathbb{Z}}
\newcommand{\Q}{\mathbb{Q}}
\newcommand{\C}{\mathbb{C}}
\newcommand{\bbH}{\mathbb{H}}


\theoremstyle{definition}
\newtheorem{ejer}{Ejercicio}
\author{Marcos Morales}
\affil{Universidad Finis Terrae}
\title{Semana 4\\}



	


\begin{document}


\begin{picture}(400, 75)
   \put(-40,15){\includegraphics[width=5cm]{UFT.png}}


\end{picture}


\begin{center}
\huge{Clase 3}
\end{center}

\begin{center}
\Large{ Marcos Morales, Camila Muñoz}
\end{center}

\begin{center}
	\large{Ecuaciones Diferenciales Ordinarias (EDO)}
\end{center}
\begin{center}
\setlength{\parindent}{0cm}\large{Clases centradas para capacitar a los estudiantes de ingeniería en Ecuaciones diferenciales ordinarias. \\}
\end{center}


\vspace{1cm}

\begin{center}
	\large{Contenidos:}
\end{center}

\setlength{\parindent}{2.95cm}- Factor integrante	\\
\phantom{abefwfewefwifefffw} - Ecuaciones homogeneas	\\
\vspace{6cm}


\begin{center}
\today
\end{center}
\begin{center}
Santiago, Chile
\end{center}


\pagestyle{empty}


\setcounter{page}{1}
\pagestyle{plain}
\setlength{\parindent}{0cm}





\newpage



\section{Factor integrante}

Supongamos que la EDO $M(x,y)dx+N(x,y)dy=0$ no es exacta, podemos asumir que existe una función $\omega(x,y)$ tal que al multiplicar la EDO por $\omega(x,y)$ la ecuación resultante

\begin{center}
$\omega(x,y) M(x,y)dx+\omega(x,y) N(x,y)dy=0$ 
\end{center}

es exacta. Para trabajar utilizaremos la notación de Euler $f_x =\displaystyle \frac{\partial f}{\partial x}$. Para que la ecuación sea exacta necesitamos que

\begin{center}
$\displaystyle \frac{\partial \left( \omega(x,y)M(x,y) \right)}{\partial y} = \frac{\partial \left( \omega(x,y)N(x,y) \right)}{\partial x}  $
\end{center}

Luego

\begin{center}
$\displaystyle \omega_yM+M_y\omega=  \omega_xN+N_x\omega $
\end{center}

\begin{center}
$\displaystyle \omega(M_y-N_x)=  \omega_xN - \omega_y M$
\end{center}

En el caso de que $\omega(x,y)=\omega(x)$, es decir $\omega$ es sólo función de $x$, entonces $\omega_y=0$, luego


\begin{center}
$\displaystyle \omega'=  \omega \frac{M_y-N_x}{N}$
\end{center}


Entonces si $\displaystyle \frac{M_y-N_x}{N}$ depende sólo de $x$, tendríamos que $\omega$ está dado por la  ecuación diferencial.

\begin{center}
$\displaystyle \frac{d\omega}{dx}=  \omega \frac{M_y-N_x}{N}$
\end{center}



 A $\omega(x,y)$ se le conoce como \textbf{Factor integrante}. \\

En el caso de que $\displaystyle \frac{N_x-M_y}{M}$ depende sólo de $y$, entonces $\omega$ está dado por la  ecuación diferencial.

\begin{center}
$\displaystyle  \frac{d\omega}{dy}=  \omega \frac{N_x-M_y}{M}$
\end{center}

(\textbf{Los siguientes casos no se vieron en clases})En el caso de que $\displaystyle \frac{N_x-M_y}{M-N}$ depende sólo $t=x+y$ como función, entonces $\omega$ está dado por la  ecuación diferencial.

\begin{center}
$\displaystyle  \frac{d\omega}{t}= \omega\frac{N_x-M_y}{M-N}$
\end{center}

Donde $\omega$ queda en función de $t=x+y$. Por último, si $\displaystyle \frac{M_x-N_y}{M+N}$ depende sólo de $s=x-y$, entonces $\omega$ está dado por la  ecuación diferencial.

\begin{center}
$\displaystyle  \frac{d\omega}{ds}= \omega\frac{M_x-N_y}{M+N}$
\end{center}

Donde $\omega$ queda en función de $s=x-y$. Existen más ocaciones en las cuales se puede calcular el factor integrante de forma eficiente, pero hay que tratar cada una de esos casos de forma individual. \\ \\

\textbf{Nota Importante: } Hay muchos más tipos de factores integrantes que van a aparecer en las ayudantías, aquí solo hemos menocionado unos cuantos. \\


\textbf{Ejercicio: }Resolver $x^2y^2dx+ (x^3y+y+3)dy=0$. \\

\textit{Pista: Encuentre el factor integrante asuminedo que dependerá sólo de $y$, en este caso la ecuación del factor integrante viene dada por  $\displaystyle  \frac{d\omega}{dy}=  \omega \frac{N_x-M_y}{M}$ } \\

\textbf{Respuesta: } Tenemos que $M(x,y)= x^2y^2$ y $N(x,y)= x^3+y+3$. La ecuación diferencial del factor integrante viene dada por


\begin{center}
$\displaystyle  \frac{d\omega}{dy}=  \omega \frac{N_x-M_y}{M} = \omega\frac{3yx^2-2yx^2}{x^2y^2}= \frac{\omega}{y}$
\end{center}


De aquí se concluye que

\begin{center}
$\displaystyle \frac{d\omega}{\omega} = \frac{dy}{y}$
\end{center}

Integrando

\begin{center}
$\ln(\omega)= \ln(y)+c$
\end{center}




\begin{center}
$\omega= ky$
\end{center}

Como el factor integrante no es único podemos tomar cualquiera, siempre elegiremos el más sencillo, en este caso sería cuando $k=1$, por lo tanto tenemos
$\omega(y)=y$. Multiplicando la EDO inicial $x^2y^2dx+ (x^3y+y+3)dy=0$ por $y$, obtenemos la nueva EDO

\begin{center}
$$x^2y^3dx+ (x^3y^2+y^2+3y)dy=0$$
\end{center}


En este caso $M(x,y)= x^2y^3$ y $N(x,y)=x^3y^2+y^2+3y$, podemos verificar que

\begin{center}
$\displaystyle M_y=N_x= 3x^2y^2$
\end{center}

Sea $f(x,y)$, la función dada por una EDO exacta, tenemos entonces que

\begin{center}
$\displaystyle \frac{\partial f}{\partial x} =M(x,y) = x^2y^3$
\end{center}

Luego

\begin{center}
$\displaystyle f(x,y) = \int x^2y^3 dx = \frac{x^3y^3}{3} +c(y)$
\end{center}

Por otra parte, sabemos que 

\begin{center}
$\displaystyle \frac{\partial f}{\partial y} = N(x,y)= x^3y^2+y^2+3y$
\end{center}

Ahora bien

\begin{center}
$\displaystyle \frac{\partial f}{\partial y} = \frac{\partial \left( \frac{x^3y^3}{3} +c(y) \right) }{\partial y} = x^3y^2+c'(y)$
\end{center}


Luego

\begin{center}
$ x^3y^2+c'(y) =  x^3y^2+y^2+3y$
\end{center}

Por lo tanto


\begin{center}
$ c'(y) =  y^2+3y$
\end{center}

Integrando



\begin{align*}
c(y) &=  \displaystyle \int y^2+3y dy \\
&=  \frac{y^3}{3} +\frac{3}{2}y^2 +c \\
\end{align*}

Reemplazando, obtenemos

\begin{center}
$\displaystyle f(x,y) =  \frac{x^3y^3}{3} +  \frac{y^3}{3} +\frac{3}{2}y^2 +c $
\end{center}

La solución de la EDO viene dada por

\begin{center}
$\displaystyle \frac{x^3y^3}{3} +  \frac{y^3}{3} +\frac{3}{2}y^2 =c $.
\end{center}

\newpage



\section{Ecuaciones Homogeneas}




\textbf{Definición: }Una función $f(x,y)$ se dice homogenea si cumple que $f(tx,ty)=t^nf(x,y)$ para todo $t \in \mathbb{R}$. \\

\textbf{Ejemplo: }Todos los polinomios en dos variables tales que cada término tenga el mismo grado como por ejemplo $x^2-xy$, $x^3+xy^2 +x^2y$ o bien $x+y$, son homogeneos. \\

Los polinomios que no tienen todos sus coeficientes del mismo grado no son homogeneos por ejemplo $x+y+2$ no es homogenea porque el término $2$ tiene grado cero, tampoco es homogenea la ecuación $x^3 +xy +y^2x$, porque el término $xy$ tiene grado 2 mientras que el término $x^3$ tiene grado 3. \\





\textbf{Definición: }Una función ecuación diferencial de la forma $M(x,y)dx +N(x,y)dy=0$ se dice homogenea si y sólo sí $M$ y $N$ son homogeneas del mismo grado. \\


\textbf{Resultado:} \textit{Dada una Ecuación diferencial homogenea $M(x,y)dx +N(x,y)dy=0$, haciendo el cambio de variable $y=zx$, la ecuación diferencial se transforma en una ecuación de variables separables}\\


Notemos que al hacer el cambio de variable $y=zx$, entonces se tiene que $\displaystyle \frac{dy}{dx}= z+x\frac{dz}{dx}$. \\


\textbf{Ejercicio: } Encontrar la solución de la ecuación diferencial $\displaystyle xy^2\frac{dy}{dx} =y^3-x^3$.\\

\textit{Respuesta: }Notemos que es una ED homogenea de grado 3, haciendo el cambio de variable $y=zx$, entonces  $\displaystyle \frac{dy}{dx}= z+x\frac{dz}{dx}$, reemplazando

\begin{center}
$\displaystyle  x(zx)^2 \left(  z+x\frac{dz}{dx} \right) =(zx)^3-x^3 $
\end{center}

\begin{center}
$\displaystyle  x^3z^2 \left(  z+x\frac{dz}{dx} \right) =x^3(z^3-1) $
\end{center}



\begin{center}
$\displaystyle  z^2 \left(  z+x\frac{dz}{dx} \right) =z^3-1 $
\end{center}

\begin{center}
$\displaystyle  z^3+xz^2\frac{dz}{dx} =z^3-1 $
\end{center}

\begin{center}
$\displaystyle  xz^2\frac{dz}{dx} =-1 $
\end{center}



\begin{center}
$\displaystyle  z^2dz =-\frac{dx}{x} $
\end{center}

Integrando


\begin{center}
$\displaystyle   \int z^2dz =- \int\frac{dx}{x} $
\end{center}

\begin{center}
$\displaystyle   \frac{z^3}{3} =-\ln(x) +c $
\end{center}

\begin{center}
$\displaystyle   \frac{\left(\frac{y}{x}\right)^3}{3} =-\ln(x) +c $
\end{center}


Que es la solución implícita del la ED inicial.









\end{document} 